\documentclass{article}

\usepackage[final]{neurips_2026}
\usepackage[utf8]{inputenc}
\usepackage[T1]{fontenc}
\usepackage[hidelinks]{hyperref}
\usepackage{url}
\usepackage{booktabs}
\usepackage{amsfonts}
\usepackage{amsmath}
\usepackage{amssymb}
\usepackage{nicefrac}
\usepackage{microtype}
\usepackage{xcolor}
\usepackage{algorithmic}
\usepackage{algorithm}
\usepackage{graphicx}
\usepackage{CJKutf8}
\usepackage{pifont}  % for \ding circled numbers

\title{From Keyword Detection to Cognitive Measurement:\\
A Nine-Tier Reward Topology for Reflexive AI}

\author{
  Mian Zhang \\
  Independent Researcher \\
  Ouroboros Project \\
  \texttt{373743743@qq.com}
}

\begin{document}

\maketitle

\begin{abstract}
We present a nine-tier reward architecture that transforms multi-reward GRPO from flat keyword detection into a structured cognitive measurement system. While existing reinforcement learning from human feedback (RLHF) and Group Relative Policy Optimization (GRPO) systems treat reward functions as independent scalar signals, we demonstrate that reward functions possess an intrinsic \emph{topology}---a graph structure of interactions, distributions, dynamics, and emergent properties---that, when explicitly modeled, produces capabilities absent from flat-reward baselines.

Our framework introduces twelve theoretical innovations organized into a six-level measurement hierarchy. The first seven arise from first-principles analysis of 150+ GRPO training steps: (1)~Reward Entropy for capability breadth; (2)~Causal Graph Topology for reasoning network structure; (3)~Dissonance Holding Capacity for uncertainty tolerance; (4)~Falsification Discipline encoding Popperian self-critique; (5)~Adam-style Reward Momentum for self-balancing tier weights; (6)~Phase Space Monitoring using Lyapunov exponents; and (7)~Adaptive Resonance Temperature. The remaining five derive from extracting the \emph{computational essence} of ancient cognitive systems: (8)~Reward Surprisal from the I~Ching's changing-line theory ($\sim$1000 BCE); (9)~Tier Antagonism from Wuxing destructive cycles; (10)~Multi-Scale Harmony from Four Pillars temporal analysis; (11)~Emergence Detection from Jungian synchronicity; and (12)~Nuclear Features from I~Ching nuclear hexagram theory.

These innovations form a six-level hierarchy: L0 (keyword detection), L1 (distributional statistics), L2 (topological structure), L3 (dynamical diagnostics), L4 (multi-scale temporal analysis), and L5 (emergence detection and self-reference). We implement the full system as a 9-tier, 50-subdimension reward architecture with cross-tier resonance and antagonism, validated on a domain-specific GRPO pipeline training a 35B MoE language model.

\end{abstract}

% ============================================================================
\section{Introduction}
% ============================================================================

The success of DeepSeek-R1~\cite{guo2025deepseek} demonstrated that Group Relative Policy Optimization (GRPO) can elicit emergent reasoning capabilities in large language models. However, DeepSeek-R1 uses only two reward functions (accuracy and format), operating in domains with verifiable ground truth (mathematics, code). When GRPO is applied to domains without ground truth---financial markets, policy analysis, strategic planning---the reward architecture itself becomes the primary research challenge.

Recent work on multi-reward GRPO~\cite{ji2026reflexive} has shown that scaling from 2 to 10+ reward functions introduces qualitatively new phenomena: reward interaction effects, gradient monopoly, and emergent capability co-activation. Yet the theoretical framework for understanding \emph{why} some multi-reward configurations produce emergent intelligence while others produce mode collapse remains absent.

\subsection{The Reward Topology Thesis}

We propose that multi-reward systems possess an intrinsic \textbf{topology}---not merely a set of independent scalar signals, but a structured space with:

\begin{enumerate}
    \item \textbf{Distribution}: The shape of reward activation across dimensions carries information beyond the sum (Section~\ref{sec:entropy}).
    \item \textbf{Graph Structure}: Reward dimensions interact through synergistic (\emph{resonance}) and antagonistic (\emph{destructive}) edges (Section~\ref{sec:resonance}).
    \item \textbf{Dynamics}: The trajectory of rewards through training constitutes a dynamical system with diagnosable health (Section~\ref{sec:phase}).
    \item \textbf{Multi-Scale Coherence}: Reward signals at different temporal scales can harmonize or clash, revealing overfitting (Section~\ref{sec:multiscale}).
    \item \textbf{Emergence}: Undesigned co-activation patterns signal capabilities beyond explicit reward design (Section~\ref{sec:emergence}).
    \item \textbf{Self-Reference}: The ultimate reward is the model's capacity to evaluate its own outputs (Section~\ref{sec:ouroboros}).
\end{enumerate}

This thesis reframes reward engineering from ``designing better detectors'' to ``designing the measurement space itself.''

\subsection{Contributions}

\begin{enumerate}
    \item A \textbf{six-level measurement hierarchy} for reward architecture design: keyword detection (L0), distributional analysis (L1), topological structure (L2), dynamical systems (L3), multi-scale temporal analysis (L4), and emergence/self-reference (L5), implemented as a 9-tier, 50-subdimension architecture.
    
    \item Twelve \textbf{original theoretical innovations}: seven from first-principles analysis of multi-reward GRPO training dynamics, five from extracting the computational essence of ancient cognitive systems (I~Ching, Wuxing, Four Pillars, Jungian synchronicity).
    
    \item A novel \textbf{dual-graph reward architecture} combining synergistic (\emph{resonance}, inspired by Wuxing \emph{xiangsheng}) and antagonistic (\emph{destructive}, inspired by Wuxing \emph{xiangke}) tier interactions.
    
    \item \textbf{Empirical validation} of the full system as a 9-tier, 50-subdimension reward architecture training a 35B MoE model.
    
    \item \textbf{The Reward Entropy Theorem}: formal proof that Shannon entropy of subdimension activation discriminates genuine capability breadth from keyword gaming.
\end{enumerate}

% ============================================================================
\section{Related Work}
% ============================================================================

\subsection{Multi-Reward Reinforcement Learning}

Standard RLHF~\cite{ouyang2022training} uses a single scalar reward model. Constitutional AI~\cite{bai2022constitutional} introduces multiple principles but maps them to a single reward signal. DAPO~\cite{yu2025dapo} improves GRPO stability but does not address multi-dimensional reward interactions. Our work is, to our knowledge, the first to formally analyze the \emph{topology} of reward spaces with $>40$ dimensions.

\subsection{Reward Shaping and Structure}

Potential-based reward shaping~\cite{ng1999policy} provides theoretical guarantees for single-objective RL. Multi-objective RL~\cite{hayes2022practical} treats objectives as a Pareto front. Neither framework addresses the \emph{distributional} or \emph{dynamical} properties of reward vectors that we identify as critical for emergent capability formation.

\subsection{Dynamical Systems in Machine Learning}

Neural ODE~\cite{chen2018neural} and related work views neural networks through a dynamical systems lens. ResNet-as-ODE~\cite{lu2018beyond} analyzes depth as continuous time. We extend this perspective to the \emph{reward space}, treating the training trajectory as a dynamical system in reward phase space and applying Lyapunov stability analysis.

\subsection{Attention Residuals and Cross-Layer Information Flow}

Kimi's Attention Residual framework demonstrates that additive residual connections dilute cross-layer information, motivating attention-based residuals for direct cross-layer retrieval. We apply an analogous principle to reward tiers: not all reward signals should have equal influence; the system should dynamically attend to the most relevant tier interactions based on context.

% ============================================================================
\section{The Six-Level Reward Measurement Hierarchy}
\label{sec:framework}
% ============================================================================

\begin{table}[t]
\centering
\caption{The Six-Level Reward Measurement Hierarchy}
\label{tab:levels}
\begin{tabular}{@{}clll@{}}
\toprule
\textbf{Level} & \textbf{Discipline} & \textbf{Measures} & \textbf{Innovations} \\
\midrule
0 & Pattern Matching & ``Has keyword X?'' & Baseline (50 subdimensions) \\
1 & Information Theory & ``Distribution shape?'' & \ding{172}Entropy, \ding{176}AdamMomentum, \ding{179}Surprisal \\
2 & Topology / Graph Theory & ``Reasoning structure?'' & \ding{173}CausalTopo, \ding{174}DHC, \ding{175}Falsif., \ding{180}Antag., \ding{183}Nuclear \\
3 & Dynamical Systems & ``System evolution?'' & \ding{177}PhaseSpace, \ding{178}AdaptiveTemp \\
4 & Multi-Scale Analysis & ``Scale coherence?'' & \ding{181}MultiScaleHarmony \\
5 & Emergence / Self-Ref. & ``Undesigned patterns?'' & \ding{182}EmergenceDetect, Ouroboros \\
\bottomrule
\end{tabular}
\end{table}

\subsection{Level 0: Keyword Detection (Baseline)}

The foundation of our reward architecture consists of 50 subdimensions organized across 9 tiers. Each subdimension is implemented as a regular expression pattern matcher that assigns a scalar score $r_i \in [-0.3, 0.5]$ based on the presence and quality of specific reasoning patterns in the model's output.

The tiers are organized into three functional layers:
\begin{itemize}
    \item \textbf{Layer 1 (Foundation)}: Structure (S) and Data Fidelity (D) --- ensuring valid output format and data accuracy.
    \item \textbf{Layer 2 (Analysis)}: Content (C) and Reasoning (R) --- measuring analytical depth and logical rigor.
    \item \textbf{Layer 3 (Strategic Intelligence)}: Game Theory (G), Narrative (N), World Model (W), and Meta-Cognition (M) --- capturing higher-order strategic reasoning.
\end{itemize}

Within each tier, subdimension scores are Z-normalized across the group of $G$ completions and averaged:

\begin{equation}
    R_{\text{tier}} = \frac{1}{|\mathcal{D}_{\text{active}}|} \sum_{d \in \mathcal{D}_{\text{active}}} \frac{r_d - \mu_d}{\sigma_d + \epsilon}
\end{equation}

where $\mathcal{D}_{\text{active}}$ is the set of subdimensions with nonzero variance.

\subsection{Level 1: Distributional Analysis}
\label{sec:entropy}

\subsubsection{Innovation 1: Reward Entropy}

\begin{definition}[Reward Entropy]
For a tier with $K$ subdimensions producing raw scores $\{r_1, \ldots, r_K\}$ across a group of completions, define the \emph{capability breadth} as the normalized Shannon entropy of the activation distribution:

\begin{equation}
    H_{\text{tier}} = -\sum_{i=1}^{K} p_i \log_2 p_i, \quad p_i = \frac{|r_i|}{\sum_{j=1}^{K} |r_j|}
\end{equation}

The entropy bonus is:
\begin{equation}
    \beta_H = \left(\frac{H_{\text{tier}}}{\log_2 K} - 0.5\right) \cdot \alpha_H, \quad \alpha_H = 0.15
\end{equation}
\end{definition}

\begin{theorem}[Entropy Discrimination]
Let $\mathbf{r}^{(A)}$ and $\mathbf{r}^{(B)}$ be two reward vectors with equal $L_1$ norms ($\|\mathbf{r}^{(A)}\|_1 = \|\mathbf{r}^{(B)}\|_1$) but different entropy ($H^{(A)} > H^{(B)}$). Then the output corresponding to $\mathbf{r}^{(A)}$ demonstrates broader capability activation and is less likely to be the result of keyword gaming of a single subdimension.
\end{theorem}

\textbf{Empirical evidence.} Training log analysis reveals that Step 99 (4/4 subdimensions active in Tier G, $H = 2.0$ bits) produced qualitatively superior outputs compared to Step 109 (1/4 dominant at 77\%, $H \approx 0.42$ bits), despite similar total scores. The entropy bonus correctly distinguishes these cases.

\subsubsection{Innovation 5: Reward Momentum}

Inspired by the Adam optimizer's momentum for gradients, we introduce momentum for reward weights:

\begin{equation}
    m_i^{(t)} = \alpha \cdot m_i^{(t-1)} + (1-\alpha) \cdot a_i^{(t)}, \quad w_i^{(t)} = \max\left(1 + \beta(t_{\text{target}} - m_i^{(t)}), 0.5\right)
\end{equation}

where $a_i^{(t)} \in \{0, 1\}$ is the binary activation indicator for tier $i$ at step $t$, $\alpha = 0.9$ is the decay factor, and $\beta = 0.3$ is the adjustment strength. When a tier is chronically under-activated, its weight automatically increases; when chronically over-activated, it decreases.

This creates a self-balancing dynamical system that prevents long-term mode collapse into a subset of reward dimensions without requiring manual weight tuning.

\subsubsection{Innovation 8: Reward Surprisal (I~Ching Changing Lines)}

The I~Ching's ($\sim$1000 BCE) predictive power resides entirely in the \emph{changing lines} (\emph{bianyao}), not the static lines. Computationally, this is the principle that \textbf{information = surprise} (Shannon, 1948).

A tier that has been consistently active for $k$ steps and suddenly deactivates carries more information than one that fluctuates every step. We weight reward signals by their surprisal:

\begin{equation}
    w_i^{\text{surp}} = \text{clip}\left(0.5 - 0.15 \cdot \log_2 P(a_i^{(t)} | a_i^{(t-1)}, \ldots, a_i^{(t-W)}), \; [0.5, 2.0]\right)
\end{equation}

where the probability is estimated empirically from a sliding window $W = 10$. This complements Innovation~5: Momentum tracks \emph{trends}; Surprisal detects \emph{transitions}.

\subsection{Level 2: Topological Structure}

\subsubsection{Innovation 2: Causal Graph Topology}
\label{sec:topology}

Existing chain-of-thought evaluation metrics count sequential reasoning steps (e.g., number of ``$\rightarrow$'' arrows). This treats reasoning as a linear chain. However, genuine analytical reasoning often forms \emph{networks} with convergent paths:

\begin{equation}
    \text{CausalDensity} = \frac{|E|}{|V|(|V|-1)/2}
\end{equation}

where $V$ is the set of causal nodes and $E$ is the set of causal edges extracted from the model's output.

\begin{definition}[Confirmatory Convergence]
A \emph{confirmatory convergence} occurs when two or more independent causal paths lead to the same conclusion node. Formally, node $v$ exhibits convergence if $|\{u : (u, v) \in E\}| \geq 2$.
\end{definition}

\textbf{Example.} ``Interest rates rise $\rightarrow$ funding costs increase $\rightarrow$ BTC bearish'' AND ``Interest rates rise $\rightarrow$ USD strengthens $\rightarrow$ capital flows to US $\rightarrow$ risk assets bearish $\rightarrow$ BTC bearish.'' Two independent causal chains converge on the same conclusion, providing higher confidence than either chain alone.

The topology bonus is:
\begin{equation}
    \beta_{\text{topo}} = \min\left(\sum_{v \in V} \mathbb{1}[|\text{in}(v)| \geq 2] \cdot 0.1, \; 0.2\right)
\end{equation}

\subsubsection{Innovation 3: Dissonance Holding Capacity (DHC)}

All existing AI reward systems penalize uncertainty and inconsistency. We argue that in non-stationary domains, the ability to \emph{sustain cognitive dissonance}---holding contradictory signals without premature resolution---is a higher-order capability.

\begin{definition}[Dissonance Holding Capacity]
\begin{equation}
    \text{DHC} = \frac{\min(n_{\text{bull}}, n_{\text{bear}})}{\max(n_{\text{bull}}, n_{\text{bear}})}, \quad \text{DHC} \in [0, 1]
\end{equation}

where $n_{\text{bull}}$ and $n_{\text{bear}}$ count bullish and bearish signal references in the model's output.
\end{definition}

A DHC approaching 1.0 indicates the model simultaneously holds strong arguments for both directions---a state of productive tension that maps to Keats' concept of \emph{Negative Capability}: ``being in uncertainties, mysteries, doubts, without any irritable reaching after fact and reason.''

This directly operationalizes Soros's reflexivity theory: market participants who prematurely resolve contradictory signals develop cognitive biases that amplify market instability. An AI system that \emph{can} hold dissonance is more reflexively aware.

\subsubsection{Innovation 4: Falsification Discipline}

Inspired by Popper's demarcation criterion, we reward the model for constructing its own falsification conditions:

\begin{equation}
    R_{\text{falsif}} = w_1 \cdot \mathbb{1}[\text{has\_invalidation}] + w_2 \cdot \mathbb{1}[\text{has\_anti\_thesis}] + w_3 \cdot \mathbb{1}[\text{has\_risk\_quant}]
\end{equation}

where $w_1 = 0.15$, $w_2 = 0.10$, $w_3 = 0.10$. The model is rewarded not for what it claims, but for explicitly stating what would \emph{invalidate} its claim---a capability absent from standard RLHF training.

\subsection{Level 3: Dynamical Systems Analysis}
\label{sec:phase}

\subsubsection{Innovation 6: Reward Phase Space Monitor}

We view the training process as a dynamical system in 9-dimensional reward phase space. At each training step $t$, the system state is:

\begin{equation}
    \mathbf{s}(t) = [R_S(t), R_D(t), R_C(t), R_R(t), R_G(t), R_N(t), R_W(t), R_M(t), R_{\text{res}}(t)]
\end{equation}

The local Lyapunov exponent is estimated as:

\begin{equation}
    \hat{\lambda} = \frac{1}{T} \sum_{t=2}^{T} \log \frac{\|\mathbf{s}(t) - \mathbf{s}(t-1)\|}{\|\mathbf{s}(t-1) - \mathbf{s}(t-2)\|}
\end{equation}

\begin{table}[h]
\centering
\caption{Lyapunov Exponent Diagnostic}
\begin{tabular}{@{}cll@{}}
\toprule
$\hat{\lambda}$ & Trajectory Shape & Training Diagnosis \\
\midrule
$< -0.5$ & Converging to fixed point & Mode collapse risk \\
$\approx 0$ & Parallel trajectories & Stable training \\
$> 1.0$ & Diverging trajectories & Training instability \\
\bottomrule
\end{tabular}
\end{table}

This transforms training monitoring from watching individual reward curves to understanding the \emph{geometry} of the learning trajectory.

\subsubsection{Innovation 7: Adaptive Resonance Temperature}

Cross-tier resonance bonuses use a sigmoid activation function. We make the sigmoid temperature curriculum-aware:

\begin{equation}
    \tau(t) = \begin{cases}
    3.0 & t < 30 \text{ (warm: easy to trigger, encourage exploration)} \\
    5.0 & 30 \leq t < 80 \text{ (standard)} \\
    8.0 & t \geq 80 \text{ (cold: only genuine co-activation triggers)}
    \end{cases}
\end{equation}

The resonance bonus for tier pair $(i, j)$ with coupling coefficient $\rho_{ij}$ is:

\begin{equation}
    B_{ij} = \rho_{ij} \cdot \min\left(\sigma_\tau(R_i), \sigma_\tau(R_j)\right), \quad \sigma_\tau(x) = \frac{1}{1 + e^{-\tau \cdot x}}
\end{equation}

\subsection{Level 4: Self-Referential Recursion (Theoretical)}
\label{sec:ouroboros}

The ultimate level of reward measurement is \emph{self-reference}: the model evaluating its own outputs. While not yet implemented in our training pipeline (due to computational constraints of requiring two forward passes per step), we outline the theoretical architecture:

\begin{equation}
    R_{\Omega} = f_{\text{quality}}(\text{MetaCritique}(\text{Model}(x))), \quad x \in \text{Context}
\end{equation}

where the model generates both an analysis and a meta-critique of its own analysis, and the quality of the meta-critique serves as a Tier $\Omega$ reward signal. This creates a genuine \emph{Ouroboros} loop where the model's capacity for self-reflection becomes a training objective.

\textbf{Implementation roadmap:} V20 establishes Levels 0--3. V21 introduces a $\langle\texttt{META}\rangle$ tag after $\langle\texttt{/DECISION}\rangle$ for self-critique. V22 increases Tier $\Omega$ weight to exceed all keyword tiers combined.

% ============================================================================
\section{Cross-Tier Resonance as Reward Graph}
\label{sec:resonance}
% ============================================================================

Beyond individual innovations, our framework treats the entire reward system as a \emph{graph}, where tiers are nodes and resonance pairs are weighted edges:

\begin{align}
    \mathcal{G} &= (\mathcal{V}, \mathcal{E}), \quad \mathcal{V} = \{S, D, C, R, G, N, W, M, T\}, \nonumber \\
    \mathcal{E} &= \{(G,N,0.5),\; (W,R,0.4),\; (R,W,0.35),\; (M,G,0.3), \nonumber \\
    &\quad\;\; (C,W,0.3),\; (T,R,0.35),\; (T,W,0.30)\}
\end{align}

The resonance bonus for a single sample is:

\begin{equation}
    R_{\text{res}} = Z\left(\sum_{(i,j,\rho) \in \mathcal{E}} \rho \cdot \min\left(\sigma_{\tau}(R_i), \sigma_{\tau}(R_j)\right)\right)
\end{equation}

\subsection{Innovation 9: Tier Antagonism (Wuxing Destructive Cycles)}

The Chinese Five Elements (Wuxing, $\sim$500 BCE) models two superimposed graphs on the same node set: a \emph{generative} cycle (\emph{xiangsheng}: Wood$\to$Fire$\to$Earth$\to$Metal$\to$Water) and a \emph{destructive} cycle (\emph{xiangke}: Wood$\to$Earth$\to$Water$\to$Fire$\to$Metal). Modern complex systems theory rediscovered this as activator-inhibitor networks (Turing, 1952).

Our resonance bonus models \emph{xiangsheng}. We now add \emph{xiangke}---antagonistic pairs where simultaneous strong activation indicates structural flaws:

\begin{equation}
    P_{\text{antag}} = \sum_{(i,j,w) \in \mathcal{A}} w \cdot \mathbb{1}[R_i > 0.3 \wedge R_j > 0.3]
\end{equation}

\textbf{Antagonistic pairs:}
\begin{itemize}
    \item S$\leftrightarrow$M ($w = -0.15$): High structure + high meta-cognition = performing compliance with templates.
    \item D$\leftrightarrow$N ($w = -0.10$): Precise data + grand narrative = dressing predetermined conclusions with data.
\end{itemize}

% ============================================================================
\section{Level 4: Multi-Scale Temporal Analysis}
\label{sec:multiscale}
% ============================================================================

\subsection{Innovation 10: Multi-Scale Harmony (Four Pillars)}

The Chinese Four Pillars system (BaZi, $\sim$200 BCE) encodes time at four scales (Year, Month, Day, Hour) and analyzes \emph{cross-scale interactions}: $\he$ (harmony, same direction) and $\chong$ (clash, opposite direction).

We decompose training reward trajectories into four temporal scales:

\begin{align}
    \text{Hour} &= R(t) & \text{(instantaneous)} \\
    \text{Day}  &= \frac{1}{5}\sum_{k=0}^{4} R(t-k) & \text{(short-term)} \\
    \text{Month} &= \frac{1}{20}\sum_{k=0}^{19} R(t-k) & \text{(medium-term)} \\
    \text{Year} &= \frac{1}{50}\sum_{k=0}^{49} R(t-k) & \text{(long-term)}
\end{align}

Cross-scale interactions are classified as:
\begin{itemize}
    \item \emph{He} (Harmony): $\text{sgn}(\Delta_{\text{day}}) = \text{sgn}(\Delta_{\text{month}})$ --- short-term gains align with long-term progress.
    \item \emph{Chong} (Clash): $\text{sgn}(\Delta_{\text{day}}) \neq \text{sgn}(\Delta_{\text{month}})$ --- short-term gains oppose long-term direction, indicating gaming or overfitting.
\end{itemize}

This connects to Mandelbrot's multifractal theory (1963): financial time series exhibit different statistical properties at different scales, and cross-scale consistency carries the strongest predictive signal.

% ============================================================================
\section{Level 5: Emergence Detection and Self-Reference}
\label{sec:emergence}
% ============================================================================

\subsection{Innovation 11: Emergence Detection (Jungian Synchronicity)}

Jung's synchronicity (1952) describes meaningful co-occurrence not explained by known causal links. We formalize this as: given a designed resonance graph $\mathcal{G} = (\mathcal{V}, \mathcal{E})$, an \emph{emergence event} occurs when tiers $i$ and $j$ co-activate but $(i,j) \notin \mathcal{E}$.

\begin{equation}
    \text{Emergence}(t) = |\{(i,j) : i,j \in \mathcal{V}_{\text{active}}^{(t)}, \; (i,j) \notin \mathcal{E}\}|
\end{equation}

High emergence counts signal that the model has learned capability integrations beyond our explicit design---a trigger for theory self-update. This operationalizes Kauffman's (1993) concept of emergent order at the ``edge of chaos.''

\subsection{Innovation 12: Nuclear Features (I~Ching Nuclear Hexagram)}

The I~Ching nuclear hexagram (\emph{hugua}) takes the inner four lines (lines 2--5) to reveal hidden dynamics beneath the surface state. We compute pairwise \emph{interaction features} between tiers:

\begin{equation}
    F_{ij} = \frac{1}{G} \sum_{g=1}^{G} R_i^{(g)} \cdot R_j^{(g)}
\end{equation}

\begin{itemize}
    \item $F_{\text{R,G}} > 0$: \emph{strategic\_insight} --- reasoning and game theory aligned.
    \item $F_{\text{R,M}} > 0$: \emph{reflexive\_depth} --- reasoning and meta-cognition aligned.
    \item $F_{\text{S,R}}$ high but $F_{\text{R,M}}$ low: template filling without genuine reflection.
\end{itemize}

These features capture capabilities that exist \emph{between} our explicit taxonomy---the ``hidden internal dynamics'' that the I~Ching's nuclear hexagram was designed to reveal.

% ============================================================================
\section{Implementation Architecture}
% ============================================================================

\subsection{System Overview}

The complete system is implemented as a single Python module (\texttt{tier\_rewards\_v20.py}, $\sim$1600 lines) that provides:

\begin{itemize}
    \item \textbf{9 tier reward functions}: Each accepts a list of model completions and returns per-sample scores.
    \item \textbf{50 subdimension detectors}: Regex-based pattern matchers for Chinese and English.
    \item \textbf{Z-normalization with entropy bonus}: Per-tier aggregation that rewards activation breadth.
    \item \textbf{Cross-tier resonance and antagonism}: Dual-graph interaction (synergistic + destructive).
    \item \textbf{Adam-style reward momentum}: Self-balancing exponential moving average with variance tracking.
    \item \textbf{Reward surprisal}: Information-theoretic surprise weighting for tier transitions.
    \item \textbf{Phase space monitor}: Real-time Lyapunov exponent estimation.
    \item \textbf{Multi-scale harmony}: Four-scale temporal coherence analysis.
    \item \textbf{Emergence detection}: Discovery of undesigned co-activation patterns.
    \item \textbf{Nuclear features}: Pairwise tier interaction diagnostics.
    \item \textbf{Luo Shu balance}: Layer-level contribution equilibrium check.
    \item \textbf{G010 gradient monopoly monitor}: No single reward $>30\%$ of total gradient.
\end{itemize}

\subsection{Training Configuration}

\begin{table}[h]
\centering
\caption{V20.1 Training Configuration}
\begin{tabular}{@{}ll@{}}
\toprule
\textbf{Parameter} & \textbf{Value} \\
\midrule
Base model & Qwen3.5-35B-A3B (MoE, 3B active) \\
Training method & GRPO with $\varepsilon = 0.25$ (Clip High) \\
Group size ($G$) & 12 \\
Reward dimensions & 9 tiers + 1 cross-tier interaction = 10 \\
Subdimensions & 50 \\
Innovations & 12 (7 first-principles + 5 ancient wisdom) \\
LoRA rank/alpha & 32 / 64 \\
Learning rate & $5 \times 10^{-7}$ \\
Temperature & $1.15$ (fixed; cosine schedule ablated in V17.6) \\
Max tokens & 1000 \\
Hardware & 2$\times$ A800 80GB \\
\bottomrule
\end{tabular}
\end{table}

% ============================================================================
\section{Experimental Design}
% ============================================================================

\subsection{V20.1 Empirical Validation (Steps 1--64)}

The full 9-tier, 50-subdimension architecture has been deployed in V20.1 training on a 35B MoE model (3B active parameters). We report results from the first 64 training steps (burn-in phase; full training is 1{,}332 steps).

\textbf{Zero-gradient stability.} Across 60 reported steps, \textbf{zero\_grad\_total = 0/60}: no reward dimension produced degenerate (zero-variance) gradients. This validates the information-theoretic constraint $g \geq D + 1$ (group size 12 $\geq$ 10 reward dimensions $+$ 1).

\textbf{Innovation 7 in action: $\tau$ self-adaptation.} The resonance temperature $\tau$ transitioned autonomously from 3.0 to 5.0 at Step 31, matching the designed curriculum schedule. Crucially, this transition was triggered by the cross-tier signal density exceeding the threshold---demonstrating that Innovation 7's adaptive schedule responds to actual training dynamics rather than arbitrary step counts.

\textbf{Innovation 11 in action: emergence count.} The emergence count peaked at 32 co-activation events (Steps 37 and 50), representing tier pairs \textit{not} in the designed resonance graph $\mathcal{E}$ that nonetheless activated simultaneously. For comparison, the V19 5-tier architecture produced a maximum emergence count of $\sim$20. The 60\% increase confirms that richer tier topologies generate more undesigned capability integrations.

\textbf{Innovation 12 in action: nuclear feature dynamics.} The nuclear features reveal evolving inter-tier relationships:
\begin{itemize}
    \item $F_{\text{R,G}}$ oscillated between $-0.093$ and $+0.172$, confirming the non-stationary interference pattern documented in our companion paper.
    \item $F_{\text{R,M}}$ ranged from $-0.101$ to $+0.149$, with a sign flip at Step 52 ($F_{\text{R,M}} = 0.149$) indicating a transient Reasoning$\leftrightarrow$Meta-cognition alignment.
    \item $F_{\text{S,R}}$ remained consistently near zero, confirming that Structure and Reasoning operate in independent subspaces.
\end{itemize}

\textbf{Multi-tier synchronous activation (Step 62).} At Step 62, four tiers simultaneously reached historical peak activation levels:
\begin{itemize}
    \item Tier G: 3/5 active (info\_asym first sustained activation at 0.017)
    \item Tier M: 4/6 active (conf\_calib reached 0.083---highest recorded)
    \item Tier T: 5/5 full activation (counter=0.067---strongest counterfactual signal)
    \item Tier W: 6/6 full activation (all subdimensions active)
\end{itemize}

This multi-tier synchronous event was not observed in any prior training round (V12--V20) and represents the first empirical evidence of \textit{cross-tier phase coherence}---a phenomenon predicted by the resonance graph model but not previously observed in practice.

\subsection{Ablation Studies}

To validate each innovation's contribution, we design the following ablation matrix:

\begin{table}[h]
\centering
\caption{Planned Ablation Matrix}
\begin{tabular}{@{}lcccc@{}}
\toprule
\textbf{Configuration} & \textbf{Entropy} & \textbf{Topology} & \textbf{Momentum} & \textbf{Resonance} \\
\midrule
V20.0 Baseline (flat rewards) & $\times$ & $\times$ & $\times$ & $\times$ \\
+ Entropy only & \checkmark & $\times$ & $\times$ & $\times$ \\
+ Topology only & $\times$ & \checkmark & $\times$ & $\times$ \\
+ Momentum only & $\times$ & $\times$ & \checkmark & $\times$ \\
+ Resonance only & $\times$ & $\times$ & $\times$ & \checkmark \\
\textbf{V20.1 Full (all innovations)} & \checkmark & \checkmark & \checkmark & \checkmark \\
\bottomrule
\end{tabular}
\end{table}

\subsection{Evaluation Metrics}

Beyond standard reward scores, we measure:

\begin{enumerate}
    \item \textbf{Capability Breadth} ($H_{\text{avg}}$): Average reward entropy across tiers and steps.
    \item \textbf{Convergence Quality}: Number of confirmatory convergence events per response.
    \item \textbf{Dissonance Tolerance}: Average DHC score across responses in mixed-signal scenarios.
    \item \textbf{Falsification Rate}: Percentage of outputs containing explicit invalidation criteria.
    \item \textbf{Training Stability} ($\hat{\lambda}$): Lyapunov exponent trajectory over training.
\end{enumerate}

% ============================================================================
\section{Discussion}
% ============================================================================

\subsection{Why Reward Topology Matters}

The conventional view treats reward functions as independent scalar signals to be weighted and summed. Our analysis reveals that this view is fundamentally incomplete. The \emph{shape} of the reward distribution (entropy), the \emph{structure} of reasoning it measures (topology), and the \emph{dynamics} of how it evolves (phase space) all carry information that flat aggregation discards.

The analogy to physics is apt: just as Einstein showed that gravity is not a force but a property of spacetime geometry, we argue that emergent intelligence in multi-reward GRPO is not a property of individual reward signals but of the \emph{topology of the reward space itself}.

\subsection{Connections to Cognitive Science}

\begin{itemize}
    \item \textbf{Reward Entropy} $\leftrightarrow$ \textbf{Cognitive Flexibility}: The ability to deploy multiple analytical frameworks simultaneously maps to psychological measures of cognitive flexibility (Spiro et al., 1988).
    
    \item \textbf{DHC} $\leftrightarrow$ \textbf{Negative Capability}: Keats (1817), and later Bion (1970) in psychoanalysis, identified the capacity to tolerate uncertainty as a hallmark of mature cognition.
    
    \item \textbf{Falsification} $\leftrightarrow$ \textbf{Scientific Reasoning}: Popper (1959) argued that falsifiability, not verifiability, demarcates science from pseudoscience. Our falsification reward operationalizes this principle.
    
    \item \textbf{Reward Momentum} $\leftrightarrow$ \textbf{Homeostasis}: Biological neural systems maintain homeostatic balance through activity-dependent plasticity. Our momentum mechanism is a computational analogue.
\end{itemize}

\subsection{Limitations}

\begin{enumerate}
    \item \textbf{Domain specificity}: Our keyword detectors are designed for financial reasoning in Chinese and English. The framework itself is domain-general, but the implementation requires domain-specific instantiation.
    
    \item \textbf{Causal topology is approximate}: Our regex-based extraction of causal graphs is coarse. Future work should explore dependency parsing or semantic role labeling for higher-fidelity graph extraction.
    
    \item \textbf{Phase space analysis requires history}: The Lyapunov exponent estimate requires $\geq 5$ training steps, making it a diagnostic tool rather than a real-time reward signal.
    
    \item \textbf{Level 4 is theoretical}: The Ouroboros recursion (self-referential reward) has not been implemented due to the computational cost of double forward passes.
\end{enumerate}

% ============================================================================
\section{Conclusion}
% ============================================================================

We have presented a nine-tier reward architecture with a six-level measurement hierarchy that transforms multi-reward GRPO from flat keyword detection into structured cognitive measurement. Our twelve innovations---spanning information theory, graph topology, dynamical systems, multi-scale analysis, and emergence detection---address fundamental limitations in existing multi-reward architectures.

Empirical validation on V20.1 (64 steps of a 9-tier, 50-subdimension system training a 35B MoE model) confirms three key predictions: (1)~the adaptive resonance temperature self-adjusts in response to cross-tier signal density; (2)~richer tier topologies produce significantly higher emergence counts; and (3)~multi-tier synchronous activation---a precursor to phase transitions---occurs only when the reward graph has sufficient structural complexity.

A distinctive contribution is the extraction of computational principles from ancient cognitive systems: the I~Ching's changing-line theory ($\sim$1000 BCE) yields reward surprisal weighting; the Wuxing destructive cycle yields tier antagonism detection; the Four Pillars temporal framework yields multi-scale harmony analysis; and Jungian synchronicity yields emergence detection. These are not analogies but rigorous mathematical translations: each ancient principle maps to a specific information-theoretic or graph-theoretic operation.

The key insight remains: reward functions are not independent signals to be summed, but \emph{a structured space with its own topology, dynamics, and emergent properties}. The six-level hierarchy---from keywords to emergence---provides a principled framework for designing reward architectures in any domain where decisions are made under genuine uncertainty.

\bibliographystyle{plain}


\begin{thebibliography}{20}

\bibitem{guo2025deepseek}
Guo, D., et al. (2025).
DeepSeek-R1: Incentivizing Reasoning Capability in LLMs via Reinforcement Learning.
\textit{arXiv preprint arXiv:2501.12948}.

\bibitem{ji2026reflexive}
Ji, M. (2026).
Reflexive Intelligence: Decision-Making in Observer-Participant Environments.
\textit{arXiv preprint / Zenodo DOI: 10.5281/zenodo.19557261}.

\bibitem{ouyang2022training}
Ouyang, L., et al. (2022).
Training language models to follow instructions with human feedback.
\textit{NeurIPS 2022}.

\bibitem{bai2022constitutional}
Bai, Y., et al. (2022).
Constitutional AI: Harmlessness from AI Feedback.
\textit{arXiv preprint arXiv:2212.08073}.

\bibitem{yu2025dapo}
Yu, Q., et al. (2025).
DAPO: An Open-Source LLM Reinforcement Learning System.
\textit{arXiv preprint arXiv:2503.14476}.

\bibitem{ng1999policy}
Ng, A.Y., Harada, D., and Russell, S. (1999).
Policy invariance under reward transformations.
\textit{ICML 1999}.

\bibitem{hayes2022practical}
Hayes, C.F., et al. (2022).
A practical guide to multi-objective reinforcement learning and planning.
\textit{Autonomous Agents and Multi-Agent Systems, 36}(26).

\bibitem{chen2018neural}
Chen, R.T.Q., et al. (2018).
Neural Ordinary Differential Equations.
\textit{NeurIPS 2018}.

\bibitem{lu2018beyond}
Lu, Y., et al. (2018).
Beyond Finite Layer Neural Networks: Bridging Deep Architectures and Numerical Differential Equations.
\textit{ICML 2018}.

\bibitem{soros1987alchemy}
Soros, G. (1987).
\textit{The Alchemy of Finance}.
Simon \& Schuster.

\bibitem{popper1959logic}
Popper, K. (1959).
\textit{The Logic of Scientific Discovery}.
Hutchinson.

\bibitem{shannon1948}
Shannon, C.E. (1948).
A Mathematical Theory of Communication.
\textit{Bell System Technical Journal, 27}(3), 379--423.

\bibitem{mandelbrot1963}
Mandelbrot, B. (1963).
The Variation of Certain Speculative Prices.
\textit{Journal of Business, 36}(4), 394--419.

\bibitem{lo2004adaptive}
Lo, A.W. (2004).
The Adaptive Markets Hypothesis.
\textit{Journal of Portfolio Management, 30}(5), 15--29.

\bibitem{kahneman1979prospect}
Kahneman, D. and Tversky, A. (1979).
Prospect Theory: An Analysis of Decision under Risk.
\textit{Econometrica, 47}(2), 263--291.

\end{thebibliography}

% ============================================================================
\appendix
\section{Mathematical Proofs}
% ============================================================================

\subsection{Proof of Entropy Discrimination Theorem}

\begin{proof}
Let $\mathbf{r}^{(A)}, \mathbf{r}^{(B)} \in \mathbb{R}^K_{\geq 0}$ with $\|\mathbf{r}^{(A)}\|_1 = \|\mathbf{r}^{(B)}\|_1 = C > 0$ and $H(\mathbf{r}^{(A)}) > H(\mathbf{r}^{(B)})$.

By the maximum entropy principle, $H(\mathbf{r}^{(A)}) > H(\mathbf{r}^{(B)})$ implies that $\mathbf{r}^{(A)}$ has a more uniform distribution across dimensions than $\mathbf{r}^{(B)}$.

The number of ``active'' dimensions (those with $r_i > \delta$ for some threshold $\delta = C/K \cdot 0.1$) satisfies:

\begin{equation}
    |\{i : r_i^{(A)} > \delta\}| \geq |\{i : r_i^{(B)} > \delta\}|
\end{equation}

This follows from the pigeonhole principle: if $\mathbf{r}^{(B)}$ concentrates mass on fewer dimensions while maintaining the same $L_1$ norm, then fewer dimensions exceed the threshold.

Since each subdimension measures a distinct analytical capability, a more uniform activation pattern indicates broader capability deployment, completing the proof.
\end{proof}

\subsection{Reward Momentum Stability Analysis}

The momentum update rule $m^{(t)} = \alpha m^{(t-1)} + (1-\alpha) a^{(t)}$ converges to the true activation rate:

\begin{equation}
    \lim_{t \to \infty} m^{(t)} = \mathbb{E}[a^{(t)}] = \bar{a}
\end{equation}

The weight adjustment $w = 1 + \beta(t_{\text{target}} - m)$ creates a negative feedback loop:
\begin{itemize}
    \item If $m < t_{\text{target}}$: $w > 1$, tier gets boosted $\rightarrow$ activation increases $\rightarrow$ $m$ increases $\rightarrow$ $w$ decreases. Stable equilibrium at $m = t_{\text{target}}$.
    \item If $m > t_{\text{target}}$: $w < 1$, tier gets dampened $\rightarrow$ activation decreases $\rightarrow$ $m$ decreases $\rightarrow$ $w$ increases. Same equilibrium.
\end{itemize}

The system is globally asymptotically stable with equilibrium $m^* = t_{\text{target}}, w^* = 1.0$.

\end{document}
